\documentclass[11pt,a4paper]{article}
\usepackage[T1]{fontenc}
\usepackage[utf8]{inputenc}
\usepackage{lmodern}
\usepackage{microtype}
\usepackage[a4paper,margin=2.25cm]{geometry}
\usepackage{longtable,booktabs,array,calc}
\usepackage{fancyvrb}
\usepackage{xcolor}
\usepackage{hyperref}
\usepackage{xurl}
\usepackage{fancyhdr}
\usepackage{enumitem}
\usepackage{parskip}
\usepackage{titlesec}
\definecolor{ddtadark}{RGB}{30,38,48}
\definecolor{ddtablue}{RGB}{38,82,126}
\hypersetup{colorlinks=true,linkcolor=ddtablue,urlcolor=ddtablue,citecolor=ddtablue}
\pagestyle{fancy}
\fancyhf{}
\lhead{Documentation-Driven Threat Analysis}
\rhead{Research work plan R17}
\cfoot{\thepage}
\setlength{\headheight}{14pt}
\setlist[itemize]{leftmargin=1.5em,itemsep=0.2em,topsep=0.25em}
\setlist[enumerate]{leftmargin=1.7em,itemsep=0.2em,topsep=0.25em}
\titleformat{\section}{\Large\bfseries\color{ddtadark}}{}{0em}{}
\titleformat{\subsection}{\large\bfseries\color{ddtadark}}{}{0em}{}
\titleformat{\subsubsection}{\normalsize\bfseries\color{ddtadark}}{}{0em}{}
\providecommand{\tightlist}{\setlength{\itemsep}{0pt}\setlength{\parskip}{0pt}}
\setlength{\emergencystretch}{4em}
\hbadness=10000
\hfuzz=20pt
\sloppy
\begin{document}
\section{DDTA research work plan after documentation-layer closure}\label{ddta-research-work-plan-after-documentation-layer-closure}

\textbf{WORK PLAN - REVISION 17}

\textbf{Status:} active BA4 execution plan after BA4-T1 projection-boundary/traceability pressure testing.

\textbf{Supersedes:} Revision 16 only for forward execution state; R1-R16 remain historical research records.

\subsection{1. Fixed prior state}\label{1-fixed-prior-state}

\begin{itemize}
\tightlist
\item
  Chapters 2-4: CLOSED / FINAL for current thesis scope.
\item
  Documentation layer: CLOSED.
\item
  BA0 responsibility/non-goals: CLOSED.
\item
  BA1 minimal BAE identity ontology: CLOSED.
\item
  BA2 relation/action vocabulary: CLOSED BY BA2-T4.
\item
  BA3 provenance/derivation/identity/lifecycle/change-revalidation mechanics: CLOSED BY BA3-T4.
\item
  \texttt{BAReferent} and \texttt{BAProposition}: ACCEPTED first-class semantic identity families.
\item
  ThreatForge is a case study/tooling instrument, not DDTA semantic authority.
\end{itemize}

\subsection{2. BA4 objective}\label{2-ba4-objective}

BA4 defines the smallest contract by which one accepted Base Analysis can support multiple derived human and methodology-oriented views without making any view project authority or allowing consumer-specific taxonomies to redefine the shared core.

BA4 remains representation-independent and does not design a general system-modeling notation.

\subsection{3. BA4-T1 - projection boundary, traceability and semantic-preservation lower-bound trial}\label{3-ba4-t1---projection-boundary-traceability-and-semantic-preservation-lower-bound-trial}

\textbf{Status: COMPLETED / PROVISIONAL PASS WITH PROJECTION-BOUNDARY LOWER-BOUND.}

T1 used one human-oriented projection and one bounded flow-oriented method projection over the same accepted facial-access BA meaning, with order/provider controls.

\subsubsection{Main dispositions}\label{main-dispositions}

\begin{verbatim}
projection derived from accepted BA                   REQUIRED
immutable projection revision                         REQUIRED
selection/coverage contract                           REQUIRED
omission == negation/absence                          REJECTED
meaning-bearing projection item trace to BA           REQUIRED
mandatory duplicate governed-source provenance        REJECTED
local item identity               REQUIRED WHEN MEANING-BEARING
role-bound multi-input trace      REQUIRED WHERE ROLES DIFFER
semantic-preserving selection/rename/group            PASS
stronger shared aggregation                           REJECTED
method-owned interpretation                           ALLOWED DOWNSTREAM
method interpretation promoted into BA                REJECTED
projection item cross-baseline lifecycle              NOT REQUIRED BY T1
BA change -> baseline-scoped projection rebuild       REQUIRED
new BAE family                                        NOT FORCED
new BA2 operator                                      NOT FORCED
BA1 / BA2 / BA3 reopen                                NOT TRIGGERED
\end{verbatim}

\subsubsection{Active provisional shape}\label{active-provisional-shape}

\begin{verbatim}
BAProjectionDescriptor                 [NOT BAE]
|- projectionKey
|- projectionRevisionKey              immutable
|- consumerPurpose
|- selectionCoverageContract
|- mappingContract
`- interpretationBoundary

BAProjectionMaterialization            [NOT BAE]
|- projectionRef
|- sourceBABaselineKey
`- item 0..*

BAProjectionItem                       [projection-local]
|- projectionItemKey
|- projectionOwnedKind
|- traceBinding 1..*
|    |- traceRoleKey
|    `- baElementRef
|- sharedSemanticRendering 0..1
`- methodOwnedInterpretation 0..1
\end{verbatim}

No one physical representation is mandated.

\subsection{4. T1 boundary invariants to preserve under further pressure}\label{4-t1-boundary-invariants-to-preserve-under-further-pressure}

\begin{enumerate}
\def\labelenumi{\arabic{enumi}.}
\tightlist
\item
  a projection does not become project authority;
\item
  shared semantic rendering must not strengthen, invert or erase material BA meaning;
\item
  method-owned interpretation stays explicitly downstream;
\item
  omission is interpreted only under an inspectable coverage contract;
\item
  source drill-down remains projection -\textgreater{} BA -\textgreater{} BA3 provenance -\textgreater{} governed source;
\item
  projection-local identity is not BAE identity;
\item
  projection rebuild follows accepted BA change rather than creating a second project lifecycle model.
\end{enumerate}

These are provisional until BA4 closure.

\subsection{5. Why BA4 remains open}\label{5-why-ba4-remains-open}

T1 used only one bounded method-oriented mapping family. It does not yet prove that the boundary survives when two methods impose incompatible taxonomies, when projection selection becomes highly lossy, or when method-owned interpretations are more strongly derived.

The current lower bound also needs pressure against diagnostic/stale views and cross-projection consistency after BA change.

\subsection{6. BA4-T2 - method-owned interpretation, coverage loss and cross-projection consistency pressure test}\label{6-ba4-t2---method-owned-interpretation-coverage-loss-and-cross-projection-consistency-pressure-test}

\textbf{NEXT / NOT EXECUTED BY THIS PLAN REVISION.}

BA4-T2 must perform only the next projection pressure step.

It must test at least:

\begin{enumerate}
\def\labelenumi{\arabic{enumi}.}
\tightlist
\item
  two bounded method-oriented projections with intentionally incompatible local taxonomies derived from the same accepted BA;
\item
  whether each method can add its own interpretation without forcing method semantics back into BA or into a shared projection ontology;
\item
  whether the revisioned projection descriptor plus role-bound BA trace is sufficient to review/rebuild method-owned interpretation;
\item
  aggressively lossy selection where omission/completeness errors are likely;
\item
  a projection intentionally exposing \texttt{DIAGNOSTIC\_UNRESOLVED}, \texttt{STALE} or review-oriented BA state without converting it into project truth;
\item
  M1-M4 plus order/WMS/provider rebuild behavior across projection families;
\item
  whether cross-projection comparison needs any new shared identity or relation beyond BA trace; and
\item
  whether a material counterexample forces only the smallest BA1/BA2/BA3/BA4 responsibility reopen.
\end{enumerate}

\subsubsection{BA4-T2 guardrails}\label{ba4-t2-guardrails}

Do not:

\begin{itemize}
\tightlist
\item
  define a complete STRIDE or STRIDE-AI ontology/schema;
\item
  define AnalysisRecord or Common Finding;
\item
  make DFD/threat-method categories BA1 types;
\item
  introduce a universal projection ontology merely to reconcile method taxonomies;
\item
  design ThreatForge classes/tables/UI;
\item
  reopen BA0-BA3 without a concrete shared-semantic counterexample;
\item
  start BA5 or BA6.
\end{itemize}

\subsubsection{BA4-T2 exit condition}\label{ba4-t2-exit-condition}

Produce a falsifiable refinement or rejection of the T1 projection lower bound and identify the next smallest BA4 pressure target. Do not close BA4 merely because two method projections can be rendered.

\subsection{7. BA5 - lexical vocabulary and optional assistance}\label{7-ba5---lexical-vocabulary-and-optional-assistance}

Only after BA4 boundaries are stable, define display labels, authoring synonyms and optional assistance. Preserve:

\begin{verbatim}
source lexical wording
        !=
stable BA semantic key
        !=
display label / authoring synonym
\end{verbatim}

BA4 projection-local labels do not preempt BA5 lexical design.

\subsection{8. BA6 - complete Base Analysis regression and closure}\label{8-ba6---complete-base-analysis-regression-and-closure}

Regress the complete BA0-BA5 design against the closed corpora and at least one structurally different holdout. BA6 may reopen only the smallest earlier responsibility on a material counterexample.

Only BA6 may close the complete Base Analysis milestone for the current thesis scope.

\subsection{9. Analysis envelope remains downstream}\label{9-analysis-envelope-remains-downstream}

Only after Base Analysis closure:

\begin{itemize}
\tightlist
\item
  A1 - AnalysisRecord / execution envelope;
\item
  A2 - common reviewed Finding boundary;
\item
  A3 - change/provenance integration;
\item
  formal methodology overlays and STRIDE/STRIDE-AI evaluations.
\end{itemize}

\subsection{10. Next authorized microstep}\label{10-next-authorized-microstep}

Only after the BA4-T1 package is reviewed, committed, pushed and remotely verified, execute:

\begin{quote}
\textbf{BA4-T2 - method-owned interpretation, coverage loss and cross-projection consistency pressure test.}
\end{quote}

Do not start BA5, BA6 or downstream analysis schemas before BA4-T2 is completed and reviewed.

\end{document}
